\setcounter{chapter}{2}
\chapter{Implementation}
\minitoc
\graphicspath{{Chapter3/figures/}}

%==============================================================================
\pagestyle{fancy}
\fancyhf{}
\fancyhead[R]{\bfseries\rightmark}
\fancyfoot[R]{\thepage}
\renewcommand{\headrulewidth}{0.5pt}
\renewcommand{\footrulewidth}{0pt}
\renewcommand{\chaptermark}[1]{\markboth{{\chaptername~\thechapter.\ #1}}{}}
\renewcommand{\sectionmark}[1]{\markright{\thechapter.\thesection~ #1}}

\begin{spacing}{1.2}
%==============================================================================
\section*{Introduction}
This chapter covers the implementation work and the bibliography.

\section{Tools and Languages}
The technical study may live in this chapter, or be performed in parallel with
the theoretical study (as the 2TUP model suggests). The goal here is to argue
for your technology choices: a short state of the art with a comparison, a
synthesis, and a justified choice of tools is expected.

\section{Application Walk-Through}
It is tempting to fill this section with every screenshot of the application
you worked so hard on -- resist. You should include screenshots, but pick
them carefully and present them as a scenario: choose one execution path,
say creating a new customer, then walk through the interfaces involved while
briefly describing the application's behaviour. Not too many images, not too
much prose -- concise, as always.

Avoid pasting raw code: nobody wants to read the contents of your classes.
You may include snippets to highlight specific features
\cite{YOUSFI2015,Latex} when truly useful. Installation steps, configuration
files and longer code excerpts belong in the appendix.

\subsection{Table Example}
You can use a tabular comparison between existing works.
See Table~\ref{tab:exple}.

\begin{table}[ht]
  \centering
  \caption{Comparison table}
  \footnotesize
  \begin{tabularx}{\linewidth}{|>{\bfseries\vspace*{\fill}}X
                              ||>{\centering\vspace*{\fill}}X
                                |>{\centering\vspace*{\fill}}X
                                |>{\centering\vspace*{\fill}}X
                                |>{\vspace*{\fill}}X<{\centering}|}
    \hline
            & \bfseries Col1 & \bfseries Col2 & \bfseries Col3 & \bfseries Col4 \\
    \hline\hline
    Row1    &      &  X  &      &      \\
    Row2    &  X   &     &      &      \\
    Row3    &  X   &  X  &  X   &  X   \\
    Row4    &  X   &     &  X   &  X   \\
    Row5    &  X   &     &  X   &  X   \\
    Row6    &  X   &     &  X   &  X   \\
    Row7    &  X   &     &  X   &      \\
    Row8    &  X   &  X  &  X   &      \\
    \hline
  \end{tabularx}
  \label{tab:exple}
\end{table}

\subsection{Code Example}
A Java snippet with syntax highlighting -- see Listing~\ref{code:java}.

\begin{lstlisting}[label=code:java,caption=Hello World in Java,language=java]
public class HelloWorld {
    // entry point
    public static void main(String[] args) {
        System.out.println("Hello, World");
    }
}
\end{lstlisting}

\section{Notes on the Bibliography}
Your bibliography must meet a number of criteria; otherwise it will get
flagged (and sometimes flagged even when you think it shouldn't -- everyone
has a personal opinion on what makes a good bibliography).\\
\begin{itemize}
  \item A bibliography in a good report should contain more books and
        articles than websites -- it's a bibliography after all. Prefer
        peer-reviewed works for definitions, rather than jumping to
        Wikipedia;
  \item Entries are usually ordered alphabetically, or by theme (and
        alphabetically within each theme);
  \item A bibliography entry generally takes this form:
  \begin{itemize}
    \item A unique identifier -- either a number \texttt{[1]} or the first
          author's name plus year, e.g.\ \texttt{[Kuntz, 1987]};
    \item If a book: author(s), title, publisher, ISBN/ISSN, publication date;
    \item If an article: author(s), title, journal or conference,
          publication date;
    \item If a website or electronic document: title, URL, date accessed;
    \item If a thesis: author, title, defending university, year, page count;
    \item Examples:
    \begin{description}
      \item $[\text{Bazin, 1992}]$ BAZIN R., REGNIER B. Antiviral treatments
            and their clinical trials. Rev.~Prat., 1992, 42, 2, p.~148--153.
      \item $[\text{Anderson, 1998}]$ ANDERSON P.JF. Checklist of criteria
            used for evaluation of metasites. [online]. University of
            Michigan, USA.\\
            \url{http://www.lib.umich.edu/megasite/critlist.html}.
            (Accessed 1998-09-11.)
    \end{description}
    \item In the body of the report, every claim taken from a reference must
          cite its identifier right after the claim. Failing to do so is
          plagiarism.
  \end{itemize}
\end{itemize}

\section*{Conclusion}
That's it.

%==============================================================================
\end{spacing}
